\section{Introduction}

BabyLM-style pretraining asks whether small language models can acquire useful linguistic generalizations from developmentally plausible data budgets rather than from web-scale corpora \citep{warstadt2023babylm,hu2024babylm,choshen2026babylm}. In this setting, the masked-language-modeling prediction budget is especially important. A model trained with a 15\% mask rate receives direct loss on only a subset of token positions, so the choice of which tokens become prediction targets may matter more than it does in large-data pretraining.

This paper asks whether small language models can pretrain more effectively when part of the masked-language-modeling budget is allocated to words that are hard for the model to predict, rather than selected uniformly at random. Put differently: can a model learn more from the same corpus and the same number of masked tokens if some masks are spent on words it currently struggles to predict? A secondary diagnostic question is which online signals are useful for identifying such words under a fixed MLM masking budget.

We study this question under a strict masking-only comparison. All systems use the same BabyLM Strict-Small data, tokenizer, masked language model architecture, optimizer, batch size, sequence length, learning-rate schedule, training length, and random seed policy. All systems also preserve the same total 15\% MLM mask rate. The only intended experimental difference is mask selection: which valid token positions are chosen for prediction.

The broad idea of adaptive masking is not new. Prior work has shown that MLM masking can be improved by choosing more useful prediction targets, changing masking schedules over training, or modifying token masking probabilities according to model predictability \citep{gu2020train,lee2022efficient,yang2023learning,zhong2023self,wilf2023difference,edman2025mask}. Our contribution is narrower and diagnostic. We ask which online signals, if any, are useful for selecting hard-to-predict words when the model, data, training schedule, and total mask rate are fixed. We use hard-to-predict operationally: a word is hard-to-predict if the current model often predicts it incorrectly or assigns it high loss during training, not because it has been externally annotated as linguistically difficult.

We compare several ways of operationalizing hard-to-predict masking. Heavy-Hitter Masking (HHM) uses repeated local entity-like words and recent high-loss tokens. CMS-Morph uses Count-Min Sketches over raw high-loss token strings and character trigrams, testing whether token-level difficulty should generalize across related surface forms. Accuracy-Morph replaces raw loss with smoothed top-1 prediction correctness, testing whether a candidate should become attractive only after it has been observed enough times and is often predicted incorrectly. Entity-Accuracy-Morph tests whether adding repeated-token pressure and late adaptive decay improves the correctness-smoothed signal.

The results are not a simple positive story. HHM worsens validation loss despite small gains on the available official diagnostics. Raw-loss CMS-Morph also fails to improve validation loss. Accuracy-Morph, however, improves matched validation loss in two seeds and produces the strongest Entity Tracking gains. This suggests that the adaptive signal matters more than simply increasing the amount of non-random masking.

This paper contributes a controlled diagnostic study of hard-to-predict word masking for small language model pretraining. Rather than changing the data, model, tokenizer, optimizer, schedule, training length, or total mask rate, we compare several online signals for choosing which words to mask. The results show that repeated-token and raw-loss signals are insufficient to improve validation loss in the current setup, while correctness-smoothed masking improves matched validation loss in two seeds and gives the strongest entity-tracking diagnostic, motivating further confirmation.
